\subsection{存储字节与分配字节}

常驻行数据只是设备分配的一部分，模型、查询工作区和临时副本另计。主码以 byte 存储时，行与列表偏移的存储为
\begin{equation}
 S_{\rm rows}=N\left(M+\left\lceil\frac{dB_r}{8}\right\rceil+8\right),
 \qquad S_{\rm offsets}=4(n_{\rm list}+1).
 \label{eq:4-row-storage}
\end{equation}
ID 和校正项贡献 $8N$ bytes。当 $M=d/8$、$B_r=8$ 时，仅编码载荷为
\[
 b_{\rm code}=d/8+d=9d/8=1.125d.
\]
这比 8 位标量量化（SQ8）多 12.5\%，因此图存储和向量编码字节数需要一起比较。

\begin{table}[!htbp]
\centering
\caption[每向量存储模型]{每向量字节数。JHQ 包含编码、ID 和校正项；SQ8 仅含向量编码；CAGRA int8 模型包含向量编码及 32 个四字节邻居。不含 offsets、模型参数、工作区及其他运行时分配。这些组件不是完整进程总量。}
\label{tab:4-bytes}
\footnotesize\setlength{\tabcolsep}{4pt}
\begin{tabular}{@{}rrrrrrrr@{}}
\toprule
$d$ & 主码 & 残差 & ID＋校正 & JHQ & SQ8 & \shortstack{CAGRA\\int8} & FP32 \\
\midrule
768 & 96 & 768 & 8 & 872 & 768 & 896 & 3,072 \\
1,024 & 128 & 1,024 & 8 & 1,160 & 1,024 & 1,152 & 4,096 \\
1,536 & 192 & 1,536 & 8 & 1,736 & 1,536 & 1,664 & 6,144 \\
3,072 & 384 & 3,072 & 8 & 3,464 & 3,072 & 3,200 & 12,288 \\
\bottomrule
\end{tabular}
\end{table}

图度数为 $d_g=32$ 时，CAGRA int8 每行使用 $d+4d_g$ bytes。Stella 的 JHQ 行载荷为 19.20\,GiB，接近 int8 CAGRA 报告的 19,656\,MiB 索引分配，但统计边界不同。改变 $\alpha$ 影响残差读取量，不改变已存储的行字节数。

基本模型存储计入旋转矩阵、主码与残差码本，以及质心和其范数：
\[
 S_{\rm model}
 =4d^2+4MKD_s+4MK_r+4n_{\rm list}(d+1)
 \quad\text{bytes}.
\]
派生数组另计。查询工作区包含两个查询矩阵、质心分数、LUT、probes 和保留候选。对预留查询容量 $b_{\rm cap}$，可选残差 LUT 增加 $4b_{\rm cap}dK_r$ bytes，默认直接评分。报告的 128-probe 配置为扫描分组预留约 93.8\,MiB，即使未使用也保留该分配（\cref{eq:5-group-workspace}）；这一大小取决于已分配容量。

输入归调用方所有。对 $n_b$ 个流水缓冲、每缓冲暂存 $a_{\rm in}$ 行，页锁定主机暂存与设备排序数组的存储为
\[
 S_{\rm pinned}=4n_ba_{\rm in}d,\qquad
 S_{\rm order}=S_{\rm inverse}=4N.
\]
逆映射是临时数组。分配临时区在 gather 前释放（\cref{eq:4-staging-rows,eq:4-assignment-scratch}），因此主机暂存、常驻数据和构建峰值具有不同生命周期。

\subsection{缓存复用与失效处理}

缓存键不含代码版本和全量数据哈希，因此代码或数据变化后需使缓存失效。参数扫描复用模型，冷构建重新拟合；仅凭召回率相近不能证明模型相同。缓存命中仍需 QR、预热、模型上传和 add。Stella 的 6.67\,s add 包含初始工作区，不含扫描校准与搜索预热，衡量的是编码调用，而非完整重启时间。

改变 $\alpha$ 无需重训。CUDA Graph 重放会复用查询缓冲，工作区在捕获前扩容。固定返回数量时，改变 batch 大小、保留数、probes、lanes、groups 或 bounds 会触发重新捕获。校准可能绕过捕获，索引拥有的缓冲在释放前保留分配。
